\section{The remaining complexity gap and a concrete research programme}
\label{main:research}

The proved results remove two local obstacles.  We can retain the ordered
attachment information needed for disk gluing without expanding a normal
surface, and we can avoid a global high-dimensional intersection iterator
when the matching cone has small coordinate factors.  A complete recognition
bound still needs a theorem connecting these operations to all knot inputs.
This section specifies that missing theorem in measurable terms and proposes
research tasks with clear success and failure conditions.

\subsection{A ledger that prevents local and global bounds from being confused}

There are at least four independent parameters in a hierarchy algorithm:
the number of states visited, the encoding length of each state, the cost of
finding a useful next surface, and the cost of verifying and updating its
geometry.  A source core of $O(T)$ handles controls only part of the last
quantity, relative to the \emph{current} source size $T$.

For example, suppose the next source has at most $30T$ exceptional handles.
Iterating this estimate through $h$ cuts permits $30^hT$ handles unless a
separate contraction or amortization theorem intervenes.  The bound does not
remain $30T$ relative to the initial triangulation.  Similarly, a sequence
of polynomial-size individual normal coordinates need not be polynomially
bounded in the original diagram if each size bound refers only to its immediate
predecessor.

The following elementary accounting proposition is useful as a design gate.
It is not asserted as a new general hierarchy theorem.

\begin{proposition}[Controlled-growth search ledger]
\label{main:ledger}
Consider a sound and complete deterministic recognition search represented
by a rooted tree of height at most $H$.  At depth $i$, suppose every state
has at most $S_i$ bits, every state offers at most $S_i^a$ successors, and
the overhead at each state, including terminal tests and successor enumeration,
is at most $S_i^b$.  In addition, the work per generated successor, including
certificate handling, is at most $S_i^b$.  Here $a,b\geq0$ are fixed constants
and $S_i\geq2$.  Then the total
work is at most
\[
 2(H+1)\exp\!\left(
       a\sum_{i=0}^{H-1}\log S_i+b\max_{0\leq i\leq H}\log S_i
                       \right)
\]
up to a polynomial initialization factor.  If $S_i\leq A^iS_0$ for a
fixed $A\geq1$, this is
\[
 \exp\!\bigl(O((H+1)\log S_0+H^2\log A+\log(H+1))\bigr).
\]
In particular, $S_0=n^{O(1)}$ and $H=O(\log n)$ give
$\exp(O((\log n)^2))=n^{O(\log n)}$.
\end{proposition}
\begin{proof}
At depth $j$, the number of nodes is at most
$\prod_{i<j}S_i^a$.  Each edge can be charged to its distinct child, so
the numbers of nodes and edges are each bounded by
$(H+1)\prod_{i<H}S_i^a$.  Every node and edge charge is at most
$\max_i S_i^b$, which gives the first estimate.  Under the growth
assumption, $\log S_i\leq\log S_0+i\log A$.  Summing this arithmetic
progression proves the second estimate and the specialization.
\end{proof}

If every state remains polynomial in $n$ independently of depth, height
$O((\log n)^2)$ instead gives $\exp(O((\log n)^3))$.  If only a fixed
multiplicative state-growth bound is known at that height, the displayed
ledger permits a fourth power of $\log n$.  This difference explains why an
explicit state-encoding bound matters even after a hierarchy-length estimate
has been obtained.  Bounds on expanded cell counts can be much worse; the
present work shows how to avoid them for a particular terminal operation.

For a fixed sector, the new candidate count
$\sum_j\binom{h_j}{d_j-1}$ gives another useful parameter test.  If
$h_j\leq n^c$ and $\max_jd_j\leq C\log n$, each block has
$n^{O(\log n)}$ candidate intersections, regardless of the sum of the
$d_j$.  That observation is a supplied-sector bound.  Enumerating exponentially
many sectors would negate it, and searching a whole tree with quasi-polynomial
branching can increase the final logarithmic exponent.  Candidate work and
branching work must be charged separately.

\subsection{Immediate geometric extension: absorb disk-base annuli}

The annular absorption proposition in the geometry section is the most direct
next implementation.  It identifies a restricted annular interface that already
has a complete topological interpretation: a disk-base bundle attached along
its entire vertical boundary is one additional two-handle.  Its raw co-core
occurrences can be read from the authenticated core-one ports of its circular
word.  Appending these rows before cancellation can remove some refusals of
the current consumer without importing a general annulus-recognition oracle.

The acceptance gate should require the whole component, not just a circle
with a suggestive profile: connected disk base, one boundary component,
product interval bundle, one full core circle, multiplicity one, and source
verification of every crossing chart.  The circle's signed walk can be checked
against the ends of the core one-handles.  For the source edge-cap face sector,
the initial endpoint's tetrahedron and face flags determine which one-handle
end is entered; both flags are needed for self-face gluings.  Taking absolute
occurrences after this geometric check supplies the saturated matrix row.

\begin{question}[Annular absorption coverage]
After this extension, how many actual source cuts become certifiable, and
which residual annuli remain?  Does absorption commute with the source leaf
schedule in a useful canonical order, or does it create a genuinely larger
positive-genus pivot search?
\end{question}

The first experiment should rerun the same 106 sources and report accepted
cuts, independent replays, explicit-oracle agreements, and each remaining
refusal category.  The existing sixty annular refusals are a denominator for
that experiment, not a claim that sixty cases will succeed after absorption.
After the existing cancellations, the added attaching rows may have no
singleton pivot, or the augmented core may not be a handlebody.

A separate source-profile diagnostic narrows the immediate target: in 38 of
the sixty annular refusals, every full-circle attachment lies on a disk base;
the other 22 have at least one nondisk-base circle.  Saved profiles resolve
45 cases, and fifteen additional ordered queries, all independently replayed,
resolve the rest.  These 38 cases satisfy the absorption eligibility condition.
They have not been reclassified by an augmented cancellation implementation.

\subsection{General annular attachments need relative topology}

An orientable interval bundle over an annulus is a solid torus.  Gluing it
to a core along one or more vertical annuli depends on how their core curves
sit in the boundary of that core.  Base Euler characteristic and an incidence
graph do not encode this information.  Even for a disk-base two-handle, the
existence of a geometric singleton is stronger than an exponent sum condition.

\begin{question}[A small relative attachment language]
Can the source words be enriched by a polynomial-size representation of a
meridian system and peripheral curves so that primitive annular attachments,
products, and selected handle cancellations are independently recognizable?
Can this be done while retaining exact source embeddings through every update?
\end{question}

For an absorbed disk-base two-handle, a certified unique intersection with
a retained geometric meridian supplies the cancellation proved by the disk-slide
lemma.  General annulus-base pieces still require relative gluing analysis.
A wider class might admit a finite sequence
of source-chart disk slides with compressed descriptions.  An arbitrary free-group
normalization should not be promoted to a geometric move unless its effect on
the boundary embedding and the remaining attaching curves is also justified.

\subsection{Transport boundary patterns through successful cancellations}

The present certificate proves that a suitable disk system exists after the
slides.  It need not construct every disk or the induced boundary
homeomorphism.  This suffices for a handlebody genus census because arbitrary
embedded disks in a handlebody boundary are legitimate gluing disks.  It is
insufficient for a next hierarchy step whose hypotheses refer to the exact
position of a boundary pattern.

\begin{question}[Compressed pattern transport]
Is there a representation of curves and trivalent patterns on the core boundary
that is closed under the meridian slides of the positive pivot lemma and has
polynomial bit growth through a successful raw epoch?  Which relative isotopies
must a checker record to preserve pattern incidence and essentiality?
\end{question}

A useful certificate should carry source flags for each pattern strand,
ordered contacts on its receiving chart, and enough local crossing data to
distinguish disjoint arcs from equal homology classes.  A directed acyclic
description of transported strands is a candidate representation, but its
size and the cost of intersection queries need proof.  The cap-two semiring
alone loses exactly the information such a transport requires.

\subsection{Handle cores with three-handles}

Ten of the saved corpus attempts stop at a core containing three-handles.
The current refusal is intentional.  An Euler calculation on the remaining
handles does not decide how a three-handle closes a spherical boundary
component, and simply deleting the corresponding row of a chain complex
does not preserve the manifold.

\begin{question}[Source-certified two/three simplification]
Can an actual source two/three cancellation or a certified spherical-boundary
reconstruction remove the three-handles while retaining the marked interfaces?
Is there a geometric subfamily where its complete search has a polynomial
certificate and producer bound?
\end{question}

A first implementation should authenticate the attaching sphere through its
source vertex-link triangles and verify the geometric cancellation condition.
One should separately investigate whether sphere components split off rather
than cancel.  This task must not substitute a homology-ball test for a ball test.

\subsection{Reduce proof volume without weakening the checker}

The large binary tests expose a practical cost that an asymptotic polynomial
statement can hide: repeated source charts and long endpoint integers occur
in several orbit certificates.  The core proof is small, but the surrounding
source and observer proofs can be tens of megabytes.

\begin{question}[Shared source proofs and persistent query systems]
Can the six certificates share a source packet, recurring interval maps, and
unchanged trace fragments through a canonical directed acyclic representation,
while independent replay still validates every query's distinct relation and
weight transport?
\end{question}

The first target is serialization and repeated parsing, not a change in the
mathematical result.  Binary integers should have one canonical exact encoding;
the existing hexadecimal transport avoids decimal conversion limits but is
not a compact binary container.  Any new schema should retain operation caps
and reject source substitution, duplicate marks, unsupported trace events,
and producer-supplied summaries without proof.  The alternative three-count
partition certificate in the interface section provides a different route
whose actual proof size can be compared with sparse weight masks.

\subsection{Measure and explain coordinate-block width}

The finest matching decomposition is exact, but its usefulness depends on the
sizes of the resulting blocks.  The bipyramid family is deliberately transparent:
it establishes a genuine asymptotic separation for the maintained iterator,
but the separation theorem makes no claim that its positive-Euler witnesses
are useful essential disks.  It does not
show that hard unknot exteriors have useful low-width witnesses.

\begin{question}[Geometric causes of small matching blocks]
Which topological or combinatorial decompositions force small matching blocks
after triangle elimination?  Which source moves preserve that property, and
which make a single block absorb most quadrilateral coordinates?
\end{question}

The diagnostic should record maximum block nullity, total nullity, projected
hyperplane counts, minimum-interaction cluster sizes, source vertex links,
and the topology of the actual returned surfaces.  Comparing only total
nullity hides the improvement.  Comparing only the smallest Q blocks hides
the possible obstruction to full standard enumeration.

\subsection{Treat Euler localization and useful-disk discovery separately}

The Euler theorem remains exact even when triangle minima couple the Q
blocks.  That is a useful asymmetry: coupling does not hurt the Euler sign
test, but it can create additional standard rays.  The explicit two-ball
counterexample proves that an implementation must keep these two APIs separate.

\begin{question}[Useful witnesses after an Euler-positive answer]
When can a block-local Euler witness be converted to a connected essential
disk with a certified boundary slope or required relative pattern condition?
When no such conversion succeeds, which further standard rays must be examined?
\end{question}

The correct experiment counts useful disks after all geometric predicates,
not merely positive values of the Euler functional.  The existing source
component and boundary-essentiality machinery is an appropriate checker for
that experiment.  A negative canonical-Euler result, a positive but inessential
disk, an incomplete search, and a verified essential disk should remain
distinct outcomes.

\subsection{The decisive missing search theorem: sector coverage}

\begin{question}[A small complete family of candidate sectors]
Can every knot exterior at the relevant stage be supplied, in quasi-polynomial
time, with a quasi-polynomial-size family of compatible sectors containing
one useful decomposing surface, with controlled local block or interaction
dimensions?  If this is false, what obstruction family shows its failure?
\end{question}

This is stronger than theorems about checking a supplied vector.  A productive
approach is to identify a progress property first and then formulate a
structural existence theorem for surfaces satisfying it.  Support-sensitive
search, matching-sign implications, linear-programming witnesses, and the new
block localization can all be candidate generators.  None alone certifies
that the family searched contains a useful surface on every input.

Possible falsification experiments should enumerate small genuine source
triangulations, certify their topology independently, and compare the proposed
restricted family against complete standard-ray searches where feasible.
Abstract matrices are valuable for algebraic obstructions but cannot establish
or refute a source-realized geometric existence claim on their own.

\subsection{Admissible patterned states and uniform source size}

The inherited core is a convenient source handle structure, not automatically
an admissible patterned hierarchy state.  The audited repository already
contains an explicit empty-cut one-tetrahedron solid torus whose inherited
characteristic-surface index violates the required admissibility condition.
Thus bounded core size does not discharge that obligation, even on the simplest
source fixture\cite{ProveIt2026}.

\begin{question}[A verified conversion to the hierarchy state model]
Can a source core, its ordered interfaces, and its transported boundary pattern
be converted to an admissible structure of uniformly controlled type with a
polynomial-size certificate?  What potential decreases under the resulting
geometric operations, and how does that potential bound both depth and restarts?
\end{question}

The recurrence in Proposition~\ref{main:ledger} should be filled with explicit
functions derived from this conversion.  It is insufficient to call every
individual state ``compressed.''  One needs a bound on the number of charts,
the length of each affine or word description, and the accumulated bit length
of ranks and markings after the allowed moves.

\subsection{Formalize the small geometric bridge}

\begin{question}[A manageable formal verification target]
Can the source chart-to-incidence lemma, the saturated positive pivot lemma,
and the disk-gluing genus theorem be formalized as a small verified kernel,
with machine-checkable hypotheses for an actual cut packet?
\end{question}

The finite arithmetic portions are especially suitable: exact integer codecs,
spanning forests, source leaf orders, permanent-one implications used only in
the meta-proof, and multigraph cycle ranks.  The topological interface is more
substantial.  It must state that the source packet represents an embedded
interval bundle and a genuine handle structure with the asserted intersection,
not merely a cell complex having the same homology.  Isolating that interface
would make both the mathematical proof and future code review more reliable.

\subsection{Evaluation and integration priorities}

The next measured comparison should keep three axes visible: source complexity,
coordinate bit length, and useful geometric coverage.  A benchmark consisting
only of enormous multiples of a known disk establishes compressed capacity but
does not measure hard surface discovery.  Conversely, a bank of tiny primitive
vectors conceals the asymptotic cost of explicit expansion.  Both are needed.

The recommended integration order is:
\begin{enumerate}
\item Preserve this additive operator and its independent source-bound checker;
keep the existing recognizer's decision contract intact.
\item Implement and audit the proved disk-base annular absorption rule.
\item Reduce repeated source/proof serialization and measure memory as well as
producer and verifier time.
\item Enable block decomposition in the matching query paths whose ray-set or
Euler contract it preserves exactly, measuring preparation and replay costs.
\item Build the relative pattern transport and useful-surface coverage theorems
needed by an actual hierarchy consumer.
\item Establish a global state and search ledger in the original diagram size
before claiming a quasi-polynomial recognizer.
\end{enumerate}

The mathematical deliverable is useful even if the last step remains difficult:
the new certificates identify precisely which topological information a
compressed cut consumer can use, and the new counterexamples identify which
apparently plausible shortcuts would lose necessary information.
